% ===================================================================
%  TABELL Nr. 13 — Den Hotel. --- De Restaurant. De Café.
%  Traduction 2026 d'apres texto/io, controlee sur texto/fr, texto/de
%  et texto/nl. Consigne : voir texto/lb/10-tabell-01.tex.
%
%  LE POURBOIRE DE L'ALINEA 9 CLOT UN FIL COMMENCE DEUX COLONNES
%  PLUS TOT, ET IL LE CLOT PAR UNE QUATRIEME SOLUTION.
%
%    l'ido et le francais : « gratifiketo », « pourboire » — un
%      compose fait chez soi ;
%    le tcheque : « spropitné » — un calque de l'allemand
%      « Trinkgeld », l'argent pour boire ;
%    le lituanien : « arbatpinigiai » — un calque du russe
%      « чаевые », l'argent pour le the ;
%    LE LUXEMBOURGEOIS : « Trinkgeld » — le mot allemand lui-meme,
%      sans une lettre de changee.
%
%  ET CELA APPREND CE QUE LE CALQUE EST. La colonne tcheque avait
%  presente le calque comme la troisieme voie entre garder le mot
%  etranger et le refaire de rien ; la colonne lituanienne avait
%  ajoute qu'il dit DE QUI l'on a copie. Il faut ajouter ceci : LE
%  CALQUE EST UN SIGNE DE DISTANCE. On calque le voisin dont on ne
%  veut pas prendre le mot ; on prend le mot du voisin dont on est
%  assez proche pour que le mot ne paraisse pas etranger. Le tcheque
%  et le lituanien ont calque parce qu'ils refusaient ; le
%  luxembourgeois n'a pas eu a refuser.
%
%  ET LE MENU DE L'ALINEA 9 EST TOUT FRANCAIS, ce qui suit la regle
%  du tableau 6 : « Krevetten », « Kaldaunen op Caener Aart »,
%  « Gigot », « Entremets », « Dessert ». La carte se lit dans la
%  langue ou l'on commande, et l'on ne commande pas dans la langue de
%  la cuisine.
% ===================================================================

%%K t13-noto-1 noto
\VUnotes{143.2mm}{%
\textit{\VUgras{(*) Fir d'Eenzelheeten iwwert d'Zëmmer, kuck d'Tabell Nr.~6.}}%
}

%%K t13-apar-1 apar
\VUsaut{3.0mm}
\VUtitre{15.0pt}{60}{TABELL Nr. 13}
\VUsaut{1.6mm}
\VUpk{10.2pt}{40}{Den Hotel. --- De Restaurant.}
\VUsaut{2.0mm}
\VUpk{10.2pt}{40}{De Café.}
\VUsaut{3.4mm}

%%K t13-01-1 p
1. --- Gëschter Owend zu Royan ukomm, hunn ech mech am „\textit{Hotel vum Ozean}“
niddergelooss, deen e Frënd mir recommandéiert hat.

%%K t13-01-2 p
Ee huet mer e schéint \VUgras{Zëmmer} \textsuperscript{(2)} um zweete
\VUgras{Stack} \textsuperscript{(1)} ginn, wou ech ganz gutt geschlof hunn (*).

%%K t13-02-1 p
2. --- Haut de Moie sinn ech aus mengem Zëmmer erausgaangen, an dat
d'\VUgras{Zëmmermeedchen} \textsuperscript{(3)} de Kaffi bruecht hat, a sinn an
d'\VUgras{Fëmmerzëmmer} \textsuperscript{(5)} gaangen, dat och als \VUgras{Liesesall}
\textsuperscript{(4)} déngt, fir d'\VUgras{Zeitungen} \textsuperscript{(6)} ze
liesen, während ech eng \VUgras{Zigarett} \textsuperscript{(8)} gefëmmt hunn. Nom
Liese sinn ech mam \VUgras{Ascenseur} \textsuperscript{(9)} erofgefuer a sinn duerch
d'Stad spadséiere gaangen.

%%K t13-03-1 p
3. --- Well ech vergiess hat, de \VUgras{Beamten} \textsuperscript{(12)} vum Hotel ze
froen, ëm wéi eng Auer d'Iesse um \VUgras{gemeinsamen Dësch} \textsuperscript{(11)}
zerwéiert gëtt, sinn ech fréi zréckkomm an hunn dat genotzt, fir am
\VUgras{Buedzëmmer} \textsuperscript{(10)} vum Hotel ze bueden. Duerno sinn ech nees
bei d'\VUgras{Keess} \textsuperscript{(15)} gaangen, fir engem vu menge Frënn ze
telefonéieren. Mä den \VUgras{Telefon} \textsuperscript{(13)} war besat; wéi
d'\VUgras{Keessiererin} \textsuperscript{(14)}, déi grad eng \VUgras{Rechnung}
\textsuperscript{(17)} an d'\VUgras{Reesendenregëster} \textsuperscript{(16)}
agedroen huet, meng Verlegenheet gesinn huet, huet si de \VUgras{Portier}
\textsuperscript{(18)} gebieden, de \VUgras{Groom} \textsuperscript{(19)} ze
schécken, fir meng Kommissioun mam \VUgras{Vëlo} \textsuperscript{(20)} ze maachen.

%%K t13-04-1 p
4. --- Ech hunn hir Merci gesot a sinn erausgaangen, fir dem \VUgras{Dréier}
\textsuperscript{(24)} (dem Handaarbechter) en Trinkgeld ze ginn, deen op senger
\VUgras{Handkar} \textsuperscript{(25)} meng \VUgras{Hutt-Këscht}
\textsuperscript{(26)}, \VUgras{meng Valis} \textsuperscript{(27)} a \VUgras{meng
Reesetäsch} \textsuperscript{(28)} vun der Gare bruecht huet. De
\VUgras{Bréifdréier} \textsuperscript{(21)}, dee grad ukomm ass, huet mer e
\VUgras{Bréif} \textsuperscript{(22)} ginn.

%%K t13-05-1 p
5. --- Ech si nach eng Zäitchen op der Schwell vun der Dier bliwwen, nieft der
\VUgras{Bréifboîte} \textsuperscript{(23)}, fir de Verkéier op der Strooss ze kucken.
De \VUgras{Kutscher} \textsuperscript{(30)} vum \VUgras{Omnibus}
\textsuperscript{(29)} huet op säi Won de \VUgras{Koffer} \textsuperscript{(31)} vun
engem \VUgras{Reesenden} \textsuperscript{(32)} gelueden, deen den
\VUgras{Hotelier} \textsuperscript{(33)} begleet huet. E jonken \VUgras{Telegrafist}
\textsuperscript{(35)} huet ouni Hascht en \VUgras{Telegramm} \textsuperscript{(34)}
fir e Client vum Hotel bruecht; en \VUgras{Tourist} \textsuperscript{(36)} mat sengem
\VUgras{Fotoapparat} \textsuperscript{(38)} iwwer der Schëller huet a säi
\VUgras{Reesféierer} \textsuperscript{(37)} gekuckt.

%%K t13-tit-1 sub
\VUsaut{4.0mm}
\VUpk{10.2pt}{40}{Mëttegiessenszäit.}
\VUsaut{3.0mm}

%%K t13-06-1 p
6. --- Virum Gitter huet e suerglose \VUgras{Commissionnaire} \textsuperscript{(39)}
tëscht sengem \VUgras{Dréierhaken} \textsuperscript{(40)} a senger
\VUgras{Cirage-Këscht} \textsuperscript{(41)} gedëmmert, während eng jonk
\VUgras{Wäschfra} \textsuperscript{(42)}, déi e \VUgras{Kuerf} \textsuperscript{(43)}
Wäsch gedroen huet, iwwert déi feierlech Mine vum \VUgras{Maître d'hôtel}
\textsuperscript{(44)} gelaacht huet, dee bei der Dier stoung.

%%K t13-07-1 p
7. --- Wéi ech de \VUgras{Kach} \textsuperscript{(45)} gesinn hunn, deen aus senger
\VUgras{Kichen} \textsuperscript{(47)} erauskoum, déi am \VUgras{Ënnergeschoss}
\textsuperscript{(48)} läit, fir e \VUgras{Botzer} \textsuperscript{(46)} eng
Kommissioun maachen ze schécken, hunn ech mech drun erënnert, datt d'Stonn vum
Mëttegiesse laanschtkënnt, a sinn an de Restaurantssall gaangen, duerch den Haff, wou
e \VUgras{Chauffeur} \textsuperscript{(50)} säin \VUgras{Auto} \textsuperscript{(49)}
gepaarkt huet, während e \VUgras{Mechaniker} \textsuperscript{(52)}, no den
Uweisunge vum \VUgras{Vëlofuerer} \textsuperscript{(53)}, en \VUgras{Petroldräirad}
\textsuperscript{(51)} gefléckt huet, dat grad en Accident hat.

%%K t13-08-1 p
8. --- Ech hunn e jonke \VUgras{Groom} \textsuperscript{(54)} gefrot, deen
\VUgras{Béiergläser} \textsuperscript{(55)} gedroen huet, ob de gemeinsamen Dësch
ënner der Veranda ass, a well hie jo gesot huet, hunn ech eng Plaz un engem vun den
Dëscher geholl an hunn mer en immens gutt Iesse ginn gelooss.

%%K t13-noto-2 noto
\VUnotes{143.2mm}{%
\textit{\VUgras{(*) Mat Hëllef vun der Platlëscht, déi am Wierderbuch steet, wäert de
Léierer dat ënnescht Menü liicht variéiere kënnen.}}%
}

%%K t13-tit-2 sub
\VUsaut{4.0mm}
\VUpk{10.2pt}{40}{En immens gutt Mëttegiessen.}
\VUsaut{3.0mm}

%%K t13-09-1 p
9. --- De Garçon huet mer d'Menü (*) vum Mëttegiessen an d'Wäilëscht bruecht. Ech
hunn eraus gesicht a bestallt als \VUgras{Virgeriicht Krevetten} mat \VUgras{Botter,}
an als Entrée \VUgras{Kaldaunen op Caener Aart.} \VUgras{Fësch} hunn ech net giess,
mä ech hunn eng Portioun \VUgras{Gigot} vu jonkem Hammel gefrot, an als
\VUgras{Geméisplat Spinat} mat Rahm. Duerno, well mer keen vun den
\VUgras{Entremets} gefall huet, hunn ech mer verschidden Desserten bréngen gelooss:
\VUgras{Kéis,} \VUgras{Uebst} a \VUgras{Kichelcher.} Ech hunn nach en immens gudde
Saint-Emilion-\VUgras{Wäin} derbäigesat, a ganz zefridde mat mengem Iessen hunn ech
den Dësch verlooss, nodeems ech dem \VUgras{Garçon} \textsuperscript{(57)} e schéint
\VUgras{Trinkgeld} \textsuperscript{(59)} ginn hunn.

%%K t13-10-1 p
10. --- Wéi de \VUgras{Gérant} \textsuperscript{(60)} mech prett fir ze goe gesinn
huet, huet hie mer proposéiert, op de Balkon ze goen, fir dat herrlecht Panorama vum
\VUgras{Ozean} \textsuperscript{(62)} ze bewonneren. Wäit ewech konnt ee
\VUgras{Fëscherbooter} \textsuperscript{(63)} gesinn, \VUgras{Segelschëffer}
\textsuperscript{(64)} an en enorme \VUgras{Paquebot} \textsuperscript{(65)}, vun deem
déi schwaarz Silhouette sech vun de wäisse \VUgras{Wollécken} \textsuperscript{(67)}
um \VUgras{Horizont} \textsuperscript{(66)} ofhieft. Eng liicht Bries huet de
\VUgras{Wimpel} \textsuperscript{(70)} flatteren gelooss, dee iwwert dem \VUgras{Schëld}
\textsuperscript{(69)} vun der \VUgras{Hal} \textsuperscript{(68)} stécht.

%%K t13-tit-3 sub
\VUsaut{3.0mm}
\VUpk{10.2pt}{40}{Vergnügen.}
\VUsaut{3.0mm}

%%K t13-11-1 p
11. --- D'Schauspill war esou interessant, datt ech beschloss hunn, muer op
d'Terrass vum Café ze goen an e \VUgras{Fernglas} \textsuperscript{(71)} oder e
\VUgras{Fernrouer} \textsuperscript{(72)} matzehuelen.

%%K t13-12-1 p
12. --- Wéi ech de Balkon verlooss hunn, sinn ech an de Café vum Hotel gaangen, fir e
\VUgras{Gedrénks} \textsuperscript{(73)} ze huelen an eng \VUgras{Billardpartie}
\textsuperscript{(75)} ze spillen. Ouni ze halen, sinn ech duerch de Café-Sall
gaangen, wou vill Gäscht \VUgras{Schach} \textsuperscript{(78)}, \VUgras{Dammespill}
\textsuperscript{(79)}, \VUgras{Dominosteng} \textsuperscript{(80)},
\VUgras{Trictrac} \textsuperscript{(81)} oder \VUgras{Kaarten}
\textsuperscript{(82)} gespillt hunn, a sinn direkt an de \VUgras{Billardsall}
\textsuperscript{(74)} eropgaangen.

%%K t13-13-1 p
13. --- Wéi d'Partie eriwwer war, sinn ech an d'Parterre erofgaangen an hunn de
Garçon ëm eng \VUgras{Enveloppe} \textsuperscript{(84)}, \VUgras{Pabeier}
\textsuperscript{(83)}, eng \VUgras{Bréifmark} \textsuperscript{(85)}, eng
\VUgras{Fieder} \textsuperscript{(87)} an \VUgras{Tënt} \textsuperscript{(86)}
gefrot, fir e Bréif ze schreiwen.

%%K t13-13-2 p
Duerno hunn ech e \VUgras{Kutscher} \textsuperscript{(92)} geruff, deem seng
\VUgras{Kutsch} \textsuperscript{(88)} virum Café stoung. Ech hunn hien no dem Präis
no de Stonnen oder fir eng Fahrt gefrot, a wéi mir eis iwwer de Präis eens gi sinn,
huet hie mer d'\VUgras{Kutschendier} \textsuperscript{(89)} opgemaach a mech
eragesat. Hien ass nees op de \VUgras{Bock} \textsuperscript{(91)} geklommen, huet
d'Päerd séier mat der \VUgras{Peitsch} \textsuperscript{(93)} ugedriwwen, a mir sinn
op e charmante Spadséiergang gefuer.

\VUsaut{6.0mm}
\VUfilet{22mm}
